\chapter{Main Work}
\label{ch:main-work}

% This is the central chapter. Replace the small algorithm below with the
% thesis's principal construction, method, or theorem sequence. State
% assumptions before claims.

\begin{assumption}[Indexed vertices]
\label{ass:indexed-vertices}
$V=\{1,\dots,n\}$, and the edges are given by a Boolean matrix
$A\in\{0,1\}^{n\times n}$.
\end{assumption}

Under \cref{ass:indexed-vertices}, the classical method of Warshall
\cite{warshall1962} updates a Boolean matrix so that, after iteration $k$, it
records paths whose intermediate vertices lie in $\{1,\dots,k\}$
(\cref{alg:warshall}).

\begin{algorithm}[htbp]
\caption{Reflexive transitive closure of an adjacency matrix}
\label{alg:warshall}
\begin{algorithmic}[1]
  \Require adjacency matrix $A\in\{0,1\}^{n\times n}$
  \Ensure reachability matrix $T$
  \State $T\gets A\lor I_n$
  \For{$k\gets 1$ to $n$}
    \For{$i\gets 1$ to $n$}
      \For{$j\gets 1$ to $n$}
        \State $T_{ij}\gets T_{ij}\lor(T_{ik}\land T_{kj})$
      \EndFor
    \EndFor
  \EndFor
  \State \Return $T$
\end{algorithmic}
\end{algorithm}

\begin{proposition}
\label{prop:warshall-correct}
When \cref{alg:warshall} terminates, $T_{ij}=1$ exactly when $j$ is reachable
from $i$.
\end{proposition}

\begin{proof}
Induct on $k$. Initially $T$ records paths of length zero or one. At iteration
$k$, the update adds precisely the paths that can be split at vertex $k$, so
after the last iteration $T$ represents the relation in
\cref{eq:transitive-closure}.
\end{proof}

The triple loop uses $O(n^3)$ Boolean operations and $O(n^2)$ memory. A thesis
should separate proved complexity bounds from measured implementation
performance.
